\documentclass[11pt,a4paper]{article}
\usepackage[margin=25mm]{geometry}
\usepackage[T1]{fontenc}
\usepackage{lmodern}
\usepackage{amsmath,amssymb,amsthm}
\usepackage{booktabs,array,longtable}
\usepackage{needspace}
\usepackage{hyperref}
\usepackage{xcolor}
\hypersetup{colorlinks=true,urlcolor=blue!55!black,linkcolor=blue!55!black,citecolor=blue!55!black,pdfauthor={Nefedov Sergey},pdftitle={A twelve-orbit cyclic 3-(31,5,1) packing}}
\newtheorem{theorem}{Theorem}
\newtheorem{lemma}{Lemma}
\newtheorem{corollary}{Corollary}
\newcommand{\ZZ}{\mathbb Z}
\title{A twelve-orbit cyclic $3$-$(31,5,1)$ packing}
\author{Nefedov Sergey\\[3pt]{\small Moscow State University, Mathematics Department}}
\date{16 October 2026 \quad Version 0.2 (working preprint)}
\begin{document}
\maketitle
\begin{abstract}
We exhibit twelve five-element subsets of $\ZZ_{31}$ whose full
translation orbits form a cyclic $3$-$(31,5,1)$ packing.  Developing these
base blocks gives $372$ distinct blocks and covers $3\,720$ distinct
triples, with none repeated. A finite certificate consisting of twelve disjoint sets of ten
triple-orbit indices is included, along with independently implemented
direct-shift and orbit-resource verifiers. A companion Lean~4 project
encodes both finite verification routes: a direct expanded-block test
using \texttt{native\_decide} and a smaller resource-table check
using kernel reduction.
Consequently the extremal number $M(31,5,3)$ of base-block orbits
satisfies $12\leq M(31,5,3)\leq13$, improving the lower bound of nine
recorded in the 25 July 2026 TheoremDB research packet. We do not
resolve the case $M=13$; no impossibility certificate is claimed.
\end{abstract}

\section{Problem and context}
Write $G=\ZZ/31\ZZ$. For $B\in\binom{G}{5}$, let
$\operatorname{Orb}(B)=\{B+t:t\in G\}$. A selection of base
blocks $B_1,\ldots,B_M$ is a cyclic $3$-packing if every triple in
$\binom{G}{3}$ belongs to at most one block in
$\bigcup_i\operatorname{Orb}(B_i)$. TheoremDB's P2632 packet
\cite{theoremdb} recorded a nine-orbit construction and the upper bound
$M\leq13$. The problem is equivalent to a length-$31$, weight-$5$
optical orthogonal code with off-peak auto-correlation and
cross-correlation at most $2$; see \cite{chung,chu} for the code setting.

\begin{theorem}\label{thm:main}
There is a cyclic $3$-$(31,5,1)$ packing with twelve base-block orbits.
In particular, $M(31,5,3)\geq12$.
\end{theorem}

\begin{table}[ht]
\centering
\caption{The twelve base blocks, each written in $[0,30]$.}
\label{tab:blocks}
\renewcommand{\arraystretch}{1.13}
\begin{tabular}{r l @{\qquad} r l}
\toprule
$i$ & $B_i$ & $i$ & $B_i$\\
\midrule
 1 & $\{0,1,2,26,29\}$ & 7 & $\{0,2,9,13,15\}$\\
 2 & $\{0,1,4,19,20\}$ & 8 & $\{0,2,10,14,25\}$\\
 3 & $\{0,1,6,14,15\}$ & 9 & $\{0,2,11,21,24\}$\\
 4 & $\{0,1,7,10,24\}$ &10 & $\{0,3,8,11,22\}$\\
 5 & $\{0,1,8,21,27\}$ &11 & $\{0,4,8,13,18\}$\\
 6 & $\{0,2,6,12,19\}$ &12 & $\{0,4,11,16,26\}$\\
\bottomrule
\end{tabular}
\end{table}

\section{A finite certificate}
For any nonempty subset $S\subseteq G$, define its translation-canonical
representative by
\[
 c(S)=\min_{t\in S}\operatorname{sort}\{(s-t)\bmod31:s\in S\},
\]
where the minimum is lexicographic. The $\binom{31}{3}=4495$ triples
partition into $145$ translation orbits, all of size $31$: no nonzero
shift can preserve a proper nonempty subset of $G$, because $31$ is prime.
Index their canonical triples in increasing lexicographic order by
$0,1,\ldots,144$. For a base block $B$, write
\[
 R(B)=\{\operatorname{index}(c(T)):T\in\textstyle{\binom B3}\}.
\]
Each row of Table~\ref{tab:certificate} specifies the indices of all ten
triple orbits developed by the corresponding block in
Table~\ref{tab:blocks}. There are \emph{exactly 120 distinct indices}
in the table (and ten per row), so the sets $R(B_i)$ are pairwise disjoint
and each has cardinality ten.

\begin{table}[ht]
\centering
\caption{Checkable certificate: ten distinct triple-orbit indices for each $B_i$.}
\label{tab:certificate}
\small
\renewcommand{\arraystretch}{1.10}
\begin{tabular}{r l}
\toprule
$i$ & indices of $R(B_i)$, in increasing order\\
\midrule
 1 & $ 0,23,24,26,27,28,50,52,53,54$\\
 2 & $ 2,10,11,14,17,18,65,66,86,87$\\
 3 & $ 4,12,13,15,16,21,97,98,112,113$\\
 4 & $ 5,8,22,64,71,72,121,123,126,134$\\
 5 & $ 6,19,25,79,92,96,109,116,129,138$\\
 6 & $30,36,43,77,84,110,111,117,128,130$\\
 7 & $33,37,39,42,44,51,89,91,120,122$\\
 8 & $34,38,49,82,88,90,114,118,133,137$\\
 9 & $35,45,48,57,59,68,127,140,141,144$\\
10 & $55,58,61,69,70,73,103,136,139,142$\\
11 & $75,76,80,81,85,94,102,106,107,135$\\
12 & $78,83,93,99,100,104,105,108,131,143$\\
\bottomrule
\end{tabular}
\end{table}

\begin{proof}[Proof of Theorem~\ref{thm:main}]
The table defines twelve distinct orbits. For a single $B_i$,
distinctness of its ten resource indices implies that no triple occurs
in two different translates of $B_i$. Pairwise disjointness of the
resource sets likewise prevents a triple from occurring in translates
of two different base blocks. Translation preserves each triple orbit.
Hence every triple is used at most once. Each chosen five-set has a
full translation orbit of size $31$. The packing therefore consists of
$12\cdot31=372$ distinct developed blocks and covers
$12\cdot31\cdot\binom53=3720$ distinct triples.
\end{proof}

\section{Upper bound and remaining gap}
A fixed pair of points lies in exactly $29$ triples. Every developed
five-block containing that pair accounts for three distinct triples
containing it. Thus the pair lies in at most $\lfloor29/3\rfloor=9$
developed blocks. Counting pair--block incidences and using $31$ blocks
per orbit gives
\[
 (31M)\binom{5}{2}\leq 9\binom{31}{2},
 \qquad\text{so}\qquad M\leq13.
\]
This recovers the previously reported bound \cite{theoremdb,bailey}.

\begin{corollary}
The maximum number of cyclic $3$-$(31,5,1)$ base-block orbits satisfies
\[\boxed{12\ \leq\ M(31,5,3)\ \leq\ 13}.\]
Only the case $M=13$ remains undecided by the present results.
\end{corollary}

For a cyclic distance $d\in\{1,\ldots,15\}$, put
$a_d(B)=\#\{\{x,y\}\subseteq B:\min(\pm(x-y)\bmod31)=d\}$.
By the same pair-incidence argument,
$\sum_i a_d(B_i)\leq9$ for each $d$. For the displayed construction the
fifteen loads are
\[
(8,7,7,9,8,8,8,8,8,8,8,9,9,8,7).
\]
These inequalities are additional cross-checks, not by themselves a
complete packing certificate.

\section{Enumeration and reproducibility}
The resource-index verifier computes $c(T)$ for every triple of a base
block. A completely separate C++ verifier instead develops all
$31$ translates and inserts every triple of every translated block into
a set; it fails immediately upon any duplicate. Both methods confirm
$372$ developed blocks, $3720$ distinct triples, and $775$ uncovered
triples out of $4495$.

The full orbit enumeration gives $5481$ orbits of five-sets, of which
$4761$ are internally admissible. The latter satisfy
$|R(B)|=10$, equivalently $|B\cap(B+d)|\leq2$ for all $d\ne0$.
The packing problem is consequently a finite $0$--$1$ set packing over
$4761$ columns and $145$ resources. The twelve-orbit certificate was
found by randomized neighborhood search after reproducing the known
nine-orbit construction; no optimality assertion is inferred from
search time, local maximality, or an incomplete MILP run.

The accompanying research bundle supplies exact input data,
source code, independent verifiers, logs, a finite certificate and a
SHA-256 digest. In particular, the sorted twelve-block plaintext
certificate has digest
\begin{center}
\ttfamily\footnotesize
bd23ca1859c55628896d6ff76cffe5b0db9291976fc0f7745195d44ce3d78e18
\end{center}
No proof of infeasibility at thirteen orbits has been produced.
Moreover, although the TheoremDB entry checked at the time of this note
still lists a lower bound of nine, publication priority requires a more
complete literature audit before any first-discovery claim.


\section{Lean 4 formalization and verification status}\label{sec:lean}
For transparency and independent reproduction, we provide a companion
Lean~4 project under \texttt{lean/} in the accompanying source archive.
The author and affiliation are those shown on the title page. The sources
pin Lean to version \texttt{v4.35.0} and use only the bundled standard
library, without a Mathlib dependency.

\paragraph{Direct expansion certificate.}
The module \texttt{lean/P2632.lean} encodes the twelve base blocks of
Table~\ref{tab:blocks} as concrete lists of residues. It constructs all
$31$ modular translates per block, sorts the resulting five-sets,
enumerates the ten index-selected triples in every developed block, and
checks that the list of $3\,720$ sorted triples has no duplicates.
It additionally checks that the $372$ developed blocks are distinct and
well formed. In the supplied source, the principal statements are
\texttt{P2632.no\_repeated\_triple} and
\texttt{P2632.exists\_twelve\_orbit\_packing}.
These larger finite decisions are expressed as
\texttt{native\_decide} proof obligations in the source.

\paragraph{Compact orbit-resource certificate.}
The module \texttt{lean/ResourceCertificate.lean} independently
recomputes the triple-orbit index of each triple from each of the twelve
base blocks and compares the results with Table~\ref{tab:certificate}.
It checks that the table contains $120$ distinct indices in the range
$0,\dots,144$. Its named result
\texttt{P2632.compact\_resource\_certificate} uses
\texttt{decide +kernel}, invoking reduction in Lean's kernel.
This module checks the finite resource data; the deduction from disjoint
translation-orbit resources to the packing condition is explained in the
mathematical proof above. We do not claim to have formalized that general
combinatorial implication as a separate Lean theorem.

\Needspace{10\baselineskip}
\paragraph{Verification methods and trust boundary.}
The two Lean routes have different trust models. In Lean~4,
\texttt{native\_decide} uses compiled native evaluation rather than a
purely kernel-reduced proof. By contrast, \texttt{decide +kernel}
reduces the smaller closed finite propositions in the kernel
\cite{leandocs}. The source files include \texttt{\#print axioms}
commands for inspection of proof dependencies. The independently
executed Python and C++ checkers provide a separate verification
route for the combinatorial witness. The associated repository also
contains a GitHub Actions workflow for reproducible builds.

To compile and inspect the proof dependencies with the pinned toolchain:
\begin{verbatim}
cd lean
lake build
lake env lean P2632.lean
lake env lean ResourceCertificate.lean
\end{verbatim}
The code, literal witness, concise resource table, and verification
instructions are distributed together. A successful Lean build would
establish the status of the accompanying formal source, but would not
resolve the global extremal question $M(31,5,3)=12$ versus $13$.

\begin{thebibliography}{9}
\bibitem{theoremdb} TheoremDB contributors,
\emph{Largest cyclic $3$-$(31,5,1)$ packing},
Problem P2632, research packet snapshot 25 July 2026.
\url{https://theoremdb.org/statements/P2632/}.
\bibitem{chung} F. R. K. Chung, J. A. Salehi and V. K. Wei,
Optical orthogonal codes: Design, analysis and applications,
\emph{IEEE Trans. Inform. Theory} \textbf{35} (1989), 595--604.
\bibitem{bailey} R. F. Bailey and A. C. Burgess,
Generalized packing designs,
\emph{Discrete Math.} \textbf{313} (2013), 1167--1190.
\bibitem{chu} W. Chu and C. J. Colbourn,
Optimal $(n,4,2)$-OOC of small orders,
\emph{Discrete Math.} \textbf{279} (2004), 163--172.
\bibitem{leandocs} Lean~4 developers,
\emph{Theorem Proving in Lean~4: Tactic Reference},
sections on \texttt{decide} and \texttt{native\_decide},
\url{https://lean-lang.org/doc/reference/latest/Tactic-Proofs/Tactic-Reference}.
\end{thebibliography}
\end{document}
